% Fuente: Control 1, P1 (pauta).
\medskip
\noindent\textbf{1.\; Modelamiento de planificación de mantenimiento y ruteo de vehículos}\par

Se propone un modelo MILP para coordinar la planificación de mantenimiento,
la asignación de equipos técnicos y el ruteo de vehículos eléctricos en un
parque eólico.

\medskip
\noindent\textbf{Conjuntos}\par

\begin{itemize}
    \item $\mathcal{T}=\{1,\dots,T\}$: conjunto de períodos de planificación.
    \item $\mathcal{I}=\{1,\dots,I\}$: conjunto de turbinas eólicas.
    \item $\mathcal{J}=\{1,\dots,J\}$: conjunto de equipos técnicos.
    \item $\delta$: depósito central.
    \item $\mathcal{V}=\mathcal{I}\cup\{\delta\}$: conjunto de nodos del problema.
\end{itemize}

\medskip
\noindent\textbf{Parámetros}\par

\begin{itemize}
    \item $\bar{C}_{it}$: costo de realizar mantenimiento a la turbina $i$ en el período $t$.
    \item $C^{tl}_{i_1 i_2}$: costo de traslado entre los nodos $i_1,i_2\in\mathcal{V}$.
    \item $\bar{o}_{ijt}$: tiempo requerido por el equipo $j$ para realizar mantenimiento a la turbina $i$ en el período $t$.
    \item $a_{ijt}\in\{0,1\}$: vale $1$ si el equipo $j$ está capacitado para mantener la turbina $i$ en el período $t$, y $0$ en caso contrario.
    \item $\tau_{i_1 i_2}$: tiempo de traslado entre los nodos $i_1,i_2\in\mathcal{V}$.
    \item $d_{i_1 i_2}$: distancia entre los nodos $i_1,i_2\in\mathcal{V}$.
    \item $\Delta_T$: duración máxima de la jornada laboral.
    \item $E$: autonomía máxima del vehículo eléctrico.
    \item $M$: número máximo de turbinas que puede atender un vehículo en un período.
\end{itemize}

\medskip
\noindent\textbf{Variables de decisión}\par

\begin{itemize}
    \item $x_{ijt}\in\{0,1\}$: vale $1$ si la turbina $i$ recibe mantenimiento por parte del equipo $j$ en el período $t$, y $0$ en caso contrario.
    \item $y_{i_1 i_2 jt}\in\{0,1\}$: vale $1$ si el equipo $j$ viaja directamente desde el nodo $i_1$ al nodo $i_2$ en el período $t$, y $0$ en caso contrario.
    \item $f_{i_1 i_2 jt}\geq 0$: variable continua de flujo usada para eliminar subtours.
\end{itemize}

\medskip
\noindent\textbf{Modelo MILP}\par

\begin{equation}
\begin{aligned}
\min \quad 
& \sum_{t\in\mathcal{T}}
  \sum_{i\in\mathcal{I}}
  \sum_{j\in\mathcal{J}}
  \bar{C}_{it}x_{ijt}
  +
  \sum_{t\in\mathcal{T}}
  \sum_{j\in\mathcal{J}}
  \sum_{i_1\in\mathcal{V}}
  \sum_{\substack{i_2\in\mathcal{V}\\ i_2\neq i_1}}
  C^{tl}_{i_1 i_2}y_{i_1 i_2 jt}.
\end{aligned}
\end{equation}

Sujeto a:

\paragraph{Asignación de mantenimiento.}

Cada turbina debe recibir mantenimiento exactamente una vez durante el horizonte:

\begin{equation}
\sum_{t\in\mathcal{T}}
\sum_{j\in\mathcal{J}}
x_{ijt}
=1,
\qquad \forall i\in\mathcal{I}.
\end{equation}

Un equipo solo puede realizar mantenimiento si está capacitado para ello:

\begin{equation}
x_{ijt}\leq a_{ijt},
\qquad \forall i\in\mathcal{I},\ \forall j\in\mathcal{J},\ \forall t\in\mathcal{T}.
\end{equation}

\paragraph{Consistencia entre mantenimiento y ruteo.}

Si una turbina recibe mantenimiento por parte de un equipo en un período, entonces dicho equipo debe llegar a la turbina:

\begin{equation}
\sum_{\substack{i_1\in\mathcal{V}\\ i_1\neq i}}
y_{i_1 ijt}
=
x_{ijt},
\qquad \forall i\in\mathcal{I},\ \forall j\in\mathcal{J},\ \forall t\in\mathcal{T}.
\end{equation}

Análogamente, si una turbina recibe mantenimiento, el equipo debe salir de ella:

\begin{equation}
\sum_{\substack{i_2\in\mathcal{V}\\ i_2\neq i}}
y_{i i_2 jt}
=
x_{ijt},
\qquad \forall i\in\mathcal{I},\ \forall j\in\mathcal{J},\ \forall t\in\mathcal{T}.
\end{equation}

Cada equipo puede iniciar a lo más una ruta desde el depósito en cada período:

\begin{equation}
\sum_{i\in\mathcal{I}}
y_{\delta ijt}
\leq 1,
\qquad \forall j\in\mathcal{J},\ \forall t\in\mathcal{T}.
\end{equation}

Si un equipo sale del depósito, debe retornar al depósito en el mismo período:

\begin{equation}
\sum_{i\in\mathcal{I}}
y_{\delta ijt}
=
\sum_{i\in\mathcal{I}}
y_{i\delta jt},
\qquad \forall j\in\mathcal{J},\ \forall t\in\mathcal{T}.
\end{equation}

\paragraph{Restricción de jornada laboral.}

La suma de los tiempos de mantenimiento y traslado no puede exceder la jornada disponible:

\begin{equation}
\sum_{i\in\mathcal{I}}
\bar{o}_{ijt}x_{ijt}
+
\sum_{i_1\in\mathcal{V}}
\sum_{\substack{i_2\in\mathcal{V}\\ i_2\neq i_1}}
\tau_{i_1 i_2}y_{i_1 i_2 jt}
\leq
\Delta_T,
\qquad \forall j\in\mathcal{J},\ \forall t\in\mathcal{T}.
\end{equation}

\paragraph{Restricción de autonomía del vehículo eléctrico.}

La distancia total recorrida por un equipo en un período no puede superar la autonomía disponible:

\begin{equation}
\sum_{i_1\in\mathcal{V}}
\sum_{\substack{i_2\in\mathcal{V}\\ i_2\neq i_1}}
d_{i_1 i_2}y_{i_1 i_2 jt}
\leq
E,
\qquad \forall j\in\mathcal{J},\ \forall t\in\mathcal{T}.
\end{equation}

\paragraph{Capacidad máxima de atención por vehículo.}

Cada vehículo puede atender a lo más $M$ turbinas en un mismo período:

\begin{equation}
\sum_{i\in\mathcal{I}}
x_{ijt}
\leq
M,
\qquad \forall j\in\mathcal{J},\ \forall t\in\mathcal{T}.
\label{ex:2:8:eq:capacidad}
\end{equation}

\paragraph{Eliminación de subtours mediante flujo.}

El flujo se interpreta como una cantidad enviada desde el depósito hacia las turbinas atendidas.
Cada turbina mantenida consume una unidad de flujo:

\begin{equation}
\sum_{\substack{i_1\in\mathcal{V}\\ i_1\neq i}}
f_{i_1 ijt}
-
\sum_{\substack{i_2\in\mathcal{V}\\ i_2\neq i}}
f_{i i_2 jt}
=
x_{ijt},
\qquad \forall i\in\mathcal{I},\ \forall j\in\mathcal{J},\ \forall t\in\mathcal{T}.
\end{equation}

El flujo solo puede circular por arcos efectivamente utilizados:

\begin{equation}
f_{i_1 i_2 jt}
\leq
M y_{i_1 i_2 jt},
\qquad 
\forall i_1,i_2\in\mathcal{V}: i_1\neq i_2,\ 
\forall j\in\mathcal{J},\ \forall t\in\mathcal{T}.
\label{ex:2:8:eq:cap-tour}
\end{equation}

Al modelar la ecuación \eqref{ex:2:8:eq:cap-tour}, la ecuación \eqref{ex:2:8:eq:capacidad} es redundante. Por lo tanto, si se modela correctamente \eqref{ex:2:8:eq:cap-tour}, no es necesario modelar \eqref{ex:2:8:eq:capacidad} (no al revés).

\paragraph{Naturaleza de las variables.}

\begin{equation}
x_{ijt}\in\{0,1\},
\qquad \forall i\in\mathcal{I},\ \forall j\in\mathcal{J},\ \forall t\in\mathcal{T}.
\end{equation}

\begin{equation}
y_{i_1 i_2 jt}\in\{0,1\},
\qquad 
\forall i_1,i_2\in\mathcal{V}: i_1\neq i_2,\ 
\forall j\in\mathcal{J},\ \forall t\in\mathcal{T}.
\end{equation}

\begin{equation}
f_{i_1 i_2 jt}\geq 0,
\qquad 
\forall i_1,i_2\in\mathcal{V}: i_1\neq i_2,\ 
\forall j\in\mathcal{J},\ \forall t\in\mathcal{T}.
\end{equation}

\medskip
\noindent\textbf{2.\; Restricción metereológica}\par

Sea $W_t\in\{0,1\}$ un parámetro que vale $1$ si las condiciones climáticas permiten
realizar mantenimiento en el período $t$, y $0$ en caso contrario.

Para incorporar esta información, basta agregar la siguiente restricción:

\begin{equation}
x_{ijt}
\leq
W_t,
\qquad \forall i\in\mathcal{I},\ \forall j\in\mathcal{J},\ \forall t\in\mathcal{T}.
\end{equation}

De esta forma, si $W_t=0$, no se puede programar mantenimiento en el período $t$.
Además, por las restricciones de consistencia entre mantenimiento y ruteo, tampoco se
activan rutas hacia turbinas en dicho período.

\medskip
\noindent\textbf{3.\; Restricción de balance de trabajo}\par

Se define la carga total de trabajo del equipo $j$ durante todo el horizonte como:

\begin{equation}
H_j
=
\sum_{t\in\mathcal{T}}
\sum_{i\in\mathcal{I}}
\bar{o}_{ijt}x_{ijt}
+
\sum_{t\in\mathcal{T}}
\sum_{i_1\in\mathcal{V}}
\sum_{\substack{i_2\in\mathcal{V}\\ i_2\neq i_1}}
\tau_{i_1 i_2}y_{i_1 i_2 jt},
\qquad \forall j\in\mathcal{J}.
\end{equation}

El requerimiento sindical exige que la diferencia entre las cargas laborales de cualquier
par de equipos no supere $\gamma$:

\begin{equation}
|H_j-H_{j'}|
\leq
\gamma,
\qquad \forall j,j'\in\mathcal{J}: j<j'.
\end{equation}

Para mantener una formulación MILP, se reemplaza el valor absoluto por dos restricciones lineales:

\begin{equation}
H_j-H_{j'}
\leq
\gamma,
\qquad \forall j,j'\in\mathcal{J}: j<j',
\end{equation}

\begin{equation}
H_{j'}-H_j
\leq
\gamma,
\qquad \forall j,j'\in\mathcal{J}: j<j'.
\end{equation}
